% Part: incompleteness
% Chapter: representability-in-q
% Section: representable-comp

\documentclass[../../../include/open-logic-section]{subfiles}

\begin{document}

\olfileid[pt]{inc}{req}{rpc}
\olsection{Funções representáveis em $\Th{Q}$ são computáveis}

Provaremos que toda função representável em~$\Th{Q}$ é computável.
Primeiro, precisamos estabelecer um lema sobre funções representáveis
em~$\Th{Q}$.

\begin{lem}\ollabel{lem:rep-q}
  Se $f(x_0,
  \dots, x_k)$ é representável em~$\Th{Q}$, existe
  !!a{formula}~$!A(x_0, \dots, x_k, y)$ tal que
  \[
  \Th{Q} \Proves !A_f(\num{n_0}, \dots, \num{n_k}, \num{m})
  \quad\text{sse}\quad m = f(n_0, \dots, n_k).
  \]
\end{lem}

\begin{proof}
  A parte ``se'' é
\olref[int]{defn:representable-fn}\olref[int]{defn:rep:a}. A parte
``somente se'' é vista da seguinte maneira: suponha que
$\Th{Q} \Proves !A_f(\num{n_0}, \dots, \num{n_k}, \num{m})$, mas
$m \neq f(n_0, \dots, n_k)$. Seja $l = f(n_0, \dots, n_k)$. Por
\olref[int]{defn:representable-fn}\olref[int]{defn:rep:a}, $\Th{Q}
\Proves !A_f(\num{n_0}, \dots, \num{n_k}, \num{l})$. Por
\olref[int]{defn:representable-fn}\olref[int]{defn:rep:b},
$\lforall[y][(!A_f(\num{n_0}, \dots, \num{n_k}, y) \lif \num{l} =
y)]$. Usando a lógica e a hipótese de que $\Th{Q} \Proves
!A_f(\num{n_0}, \dots, \num{n_k}, \num{m})$, obtemos que $\Th{Q}
\Proves \eq[\num{l}][\num{m}]$. Por outro lado, por
\olref[bre]{lem:q-proves-neq}, $\Th{Q} \Proves
\eq/[\num{l}][\num{m}]$. Logo, $\Th{Q}$ é inconsistente. Mas isso é
impossível, pois $\Th{Q}$ é satisfeita pelo modelo padrão (veja
\olref[int][def]{def:standard-model}), $\Sat{N}{\Th{Q}}$, e
sistemas satisfazíveis são sempre consistentes pelo Teorema da Correção
(\tagrefs{prfAX/{fol:axd:sou:cor:consistency-soundness},
prfSC/{fol:seq:sou:cor:consistency-soundness},
prfND/{fol:ntd:sou:cor:consistency-soundness},
prfTab/{fol:tab:sou:cor:consistency-soundness}}).
\end{proof}

\begin{lem}
Toda função representável em $\Th{Q}$ é computável.
\end{lem}

\begin{proof}
Comecemos pela ideia intuitiva de por que isso é verdade. Para
computar~$f$, fazemos o seguinte. Listamos todas as !!{derivation}s
possíveis~$\delta$ na linguagem da aritmética. Isso pode ser feito
mecanicamente. Para cada uma, verificamos se é !!a{derivation} de
!!a{formula} da forma~$!A_f(\num{n_0}, \dots,
\num{n_k}, \num{m})$ (a !!{formula} que representa $f$ em~$\Th{Q}$,
de \olref{lem:rep-q}). Se for, $m = f(n_0, \dots, n_k)$ por
\olref{lem:rep-q}, e encontramos o valor de~$f$. A busca termina porque
$\Th{Q} \Proves !A_f(\num{n_0}, \dots, \num{n_k},
\num{f(n_0, \dots, n_k)})$; portanto, acabamos encontrando um $\delta$
do tipo correto.

Isso ainda não é totalmente preciso, pois nosso procedimento opera
sobre !!{derivation}s e !!{formula}s, e não apenas sobre números, e
não explicamos exatamente por que é mecanicamente possível ``listar
todas as !!{derivation}s possíveis''. Mas, como vimos, é possível
codificar termos, !!{formula}s e !!{derivation}s por números de G\"odel.
Também introduzimos um modelo preciso de computação, as funções
recursivas gerais. E vimos que a relação $\Prf[\Th{Q}](d,y)$, que vale
sse $d$ é o número de G\"odel de !!a{derivation} da !!{formula} com
número de G\"odel~$y$ a partir dos axiomas de~$\Th{Q}$, é recursiva
(primitiva). Outras funções recursivas primitivas de que precisaremos
são $\fn{num}$ (\olref[art][trm]{prop:num-primrec}) e $\fn{Subst}$
(\olref[art][sub]{prop:subst-primrec}). A partir delas, é possível
definir~$f$ por minimização; portanto, $f$ é recursiva.

Primeiro, defina
\begin{multline*}
  A(n_0, \dots, n_k, m) = \\
  \fn{Subst}(\fn{Subst}(\dots\fn{Subst}(\Gn{!A_f}, \fn{num}(n_0), \Gn{x_0}),\\ \dots),
  \fn{num}(n_k),  \Gn{x_k}), \fn{num}(m), \Gn{y})
\end{multline*}
Isso parece complicado, mas é apenas a função $A(n_0, \dots, n_k,
m) = \Gn{!A_f(\num{n_0}, \dots, \num{n_k}, \num{m})}$.

Considere agora a relação~$R(n_0, \dots, n_k, s)$, que vale se
$(s)_0$ é o número de G\"odel de !!a{derivation} a partir de~$\Th{Q}$ de
$!A_f(\num{n_0}, \dots, \num{n_k}, \num{(s)_1})$:
\[
R(n_0, \dots, n_k, s) \quad\text{sse}\quad \Prf[\Th{Q}]((s)_0, A(n_0,
\dots, n_k, (s)_1))
\]
Se encontrarmos um~$s$ tal que $R(n_0, \dots, n_k, s)$ vale, teremos
encontrado um par de números---$(s)_0$ e~$(s)_1$---tal que $(s)_0$ é o
número de G\"odel de !!a{derivation} de~$A_f(\num{n_0}, \dots,
\num{n_k}, (s)_1)$. Assim, procurar~$s$ equivale a procurar o par $d$
e $m$ na prova informal. E uma função computável que ``procure'' tal
$s$ pode ser definida por minimização regular. Observe que $R$ é
regular: para todos $n_0$, \dots, $n_k$, existe !!a{derivation}~$\delta$
de $\Th{Q} \Proves !A_f(\num{n_0}, \dots,
\num{n_k}, \num{f(n_0, \dots, n_k)})$; logo,
$R(n_0, \dots, n_k, s)$ vale para
$s = \tuple{\Gn{\delta}, f(n_0, \dots, n_k)}$. Portanto, podemos
escrever $f$ como
\[
f(n_0,\dots,n_{k}) = (\umin{s}{R(n_0, \dots, n_k, s)})_1.
\]
\end{proof}

\end{document}
